% ================================================================ CHAPTER 3
% OBE 2.1 structure: 3.1 Results and Analysis, 3.2 Accuracy and
% Verification, 3.3 Impact Analysis.
\chapter{Results and Analysis}
\label{ch:results}

\section{Results and Analysis}
\label{sec:results}

\todo{This section is populated as experiments complete. The table skeletons
below define the empirical evidence for each research question.}

The proposed solution answers the research questions through four families of
evidence: in-domain attribution (RQ1), cross-dataset and cross-generator
generalization (RQ3, RQ6), temporal localization (RQ4), and module ablations
(RQ2, RQ6). The basic functionality of the proposed solution---independent
per-modality verdicts with quadrant attribution and manipulation
timelines---makes use of state-of-the-art self-supervised encoders and the
QACP pretraining strategy introduced in Chapter~\ref{ch:methodology}.

\begin{table}[h]
\centering
\caption{Per-modality and quadrant results on FakeAVCeleb (subject-disjoint test split). \todo{fill}}
\label{tab:indomain}
\small
\begin{tabular}{lccccc}
\toprule
\textbf{Method} & \textbf{Video AUC} & \textbf{Audio AUC} & \textbf{Quadrant Acc.} & \textbf{Macro-F1} & \textbf{EER} \\
\midrule
Xception (video-only)     & -- & -- & -- & -- & -- \\
AASIST (audio-only)       & -- & -- & -- & -- & -- \\
MDS                       & -- & -- & -- & -- & -- \\
AVoiD-DF                  & -- & -- & -- & -- & -- \\
\textbf{Proposed (ours)}  & -- & -- & -- & -- & -- \\
\bottomrule
\end{tabular}
\end{table}

\begin{table}[h]
\centering
\caption{Cross-dataset AUC: train on FakeAVCeleb + AV-Deepfake1M, test unseen. \todo{fill}}
\label{tab:crossdataset}
\small
\begin{tabular}{lccc}
\toprule
\textbf{Method} & \textbf{DFDC} & \textbf{KoDF} & \textbf{DF-TIMIT} \\
\midrule
Best unimodal baseline      & -- & -- & -- \\
Best multimodal baseline    & -- & -- & -- \\
Proposed w/o QACP           & -- & -- & -- \\
\textbf{Proposed + QACP}    & -- & -- & -- \\
\bottomrule
\end{tabular}
\end{table}

\begin{table}[h]
\centering
\caption{Ablation study. Each row reports in-domain and cross-dataset AUC. \todo{fill}}
\label{tab:ablation}
\small
\begin{tabular}{lcc}
\toprule
\textbf{Configuration} & \textbf{In-domain AUC} & \textbf{Cross-dataset AUC} \\
\midrule
Full framework (+ QACP)          & -- & -- \\
$-$ QACP pretraining             & -- & -- \\
$-$ MISMATCH pseudo-class        & -- & -- \\
$-$ synchronization module       & -- & -- \\
$-$ disentanglement loss         & -- & -- \\
early fusion                     & -- & -- \\
late fusion                      & -- & -- \\
single-task (binary only)        & -- & -- \\
$-$ localization head            & -- & -- \\
\bottomrule
\end{tabular}
\end{table}

\todo{Add temporal localization table (AP@IoU on AV-Deepfake1M and LAV-DF vs.\
BA-TFD+) and a discussion connecting each result to its RQ.}

% ------------------------------------------------------------------
% Auto-generated results figures. Produced by:
%   python -m src.eval.figures --results results/ --out report/figures/generated
% They appear automatically once the result JSONs exist (HANDOFF.md §5, §7).
% ------------------------------------------------------------------
\IfFileExists{figures/generated/results_roc.pdf}{%
\begin{figure}[h]
  \centering
  \includegraphics[width=\textwidth]{generated/results_roc}
  \caption{ROC curves of the proposed framework and baselines, per modality.}
  \label{fig:roc}
\end{figure}}{}

\IfFileExists{figures/generated/results_confusion.pdf}{%
\begin{figure}[h]
  \centering
  \includegraphics[width=0.55\textwidth]{generated/results_confusion}
  \caption{Quadrant confusion matrix on the FakeAVCeleb test split.}
  \label{fig:confusion}
\end{figure}}{}

\IfFileExists{figures/generated/results_ablation.pdf}{%
\begin{figure}[h]
  \centering
  \includegraphics[width=0.9\textwidth]{generated/results_ablation}
  \caption{Ablation study: in-domain versus cross-dataset AUC per configuration.}
  \label{fig:ablation}
\end{figure}}{}

\IfFileExists{figures/generated/results_localization.pdf}{%
\begin{figure}[h]
  \centering
  \includegraphics[width=0.9\textwidth]{generated/results_localization}
  \caption{Qualitative temporal localization: per-frame manipulation probability
  against ground-truth manipulated intervals (shaded).}
  \label{fig:localization}
\end{figure}}{}

\IfFileExists{figures/generated/results_reliability.pdf}{%
\begin{figure}[h]
  \centering
  \includegraphics[width=\textwidth]{generated/results_reliability}
  \caption{Reliability diagrams per modality with expected calibration error.}
  \label{fig:reliability}
\end{figure}}{}

\IfFileExists{figures/generated/results_robustness.pdf}{%
\begin{figure}[h]
  \centering
  \includegraphics[width=\textwidth]{generated/results_robustness}
  \caption{Robustness: AUC under video degradations (blur, resolution loss,
  quantization) and audio noise (SNR).}
  \label{fig:robustness}
\end{figure}}{}

\section{Accuracy and Verification}
\label{sec:verification}

\todo{Fill after experiments.} The accuracy of the findings is verified through
four mechanisms, each establishing a direct connection between the collected
research data and the research objective. First, \emph{statistical
verification}: every headline metric is reported as the mean and standard
deviation over three seeds with bootstrap 95\% confidence intervals, and AUC
differences against baselines are tested with DeLong tests under
Holm--Bonferroni correction. Second, \emph{calibration}: expected calibration
error and reliability diagrams verify that reported confidences are
trustworthy, which is essential for forensic use. Third, \emph{robustness}:
detection quality is measured under H.264 re-compression, audio noise, and
resolution loss to identify potential points of failure, and the observed
degradation profile defines the operating envelope of the system. Fourth,
\emph{external validity}: cross-dataset and leave-one-generator-out protocols
verify that the solution is resilient to distribution shift rather than fitted
to one benchmark. Trade-offs between the full research model and the
lightweight deployed variant (accuracy versus latency and hosting cost) are
presented to reflect realistic product requirements.

\section{Impact Analysis}
\label{sec:impact}

\textbf{Societal impact.} Reliable per-modality attribution strengthens the
information ecosystem: fact-checkers can state precisely which component of a
viral clip was fabricated, platforms can apply proportionate moderation, and
citizens gain a free verification tool. Conversely, detector outputs could be
misused to dismiss authentic evidence (``liar's dividend''); the system
mitigates this by reporting calibrated probabilities with explicit uncertainty
rather than categorical verdicts.

\textbf{Legal and safety issues.} Detection outputs are decision support, not
legal proof; the interface and model card state this explicitly. The
framework's localization output supports human expert review, which remains
the appropriate standard in legal contexts.

\textbf{Cultural and fairness issues.} Because training corpora
over-represent some demographics and languages, the fairness audit
(Section~\ref{sec:verification}) quantifies subgroup gaps, and balanced
sampling reduces them; residual gaps are documented so downstream users
understand for whom the detector is less reliable.

\textbf{Sustainability.} The cached-feature training regime and the
lightweight deployed variant minimize energy consumption relative to
end-to-end training and serving of large models. The QACP construction is
self-refreshing---new pseudo-forgeries can be synthesized from pristine data
as generators evolve---so the solution remains maintainable without repeated
large-scale data collection, benefiting stakeholders from platform operators
to end users.
